\begin{proposition}[可观测最小多项式的非线性阶塌缩：\texorpdfstring{$q=9,10,11,13$}{q=9,10,11,13} 首次出现固定亏格（可审计）]\label{prop:pom-observable-minpoly-order-collapse-q9-13}
沿用定理 \ref{thm:pom-moment-resonance} 的记号，令 \(P_q(\lambda)\in\ZZ[\lambda]\) 为序列 \(m\mapsto S_q(m)\) 的最小整系数递推特征多项式（亦即可观测最小湮灭多项式；注意本文用 \(\Pi_q\) 表示 Perron 投影，故不采用 \(\Pi_q(\lambda)\) 记号以免混淆）。
则在 \(q\in\{9,10,11,13\}\) 上，\(P_q(\lambda)\) 为不可约整系数多项式，其次数分别为
\[
\deg P_9=7,\qquad \deg P_{10}=9,\qquad \deg P_{11}=9,\qquad \deg P_{13}=11,
\]
并且相对差分模态截断的“名义维数” \(2\lfloor q/2\rfloor+1\) 均出现固定亏格
\[
2\Bigl\lfloor\frac{q}{2}\Bigr\rfloor+1-\deg P_q=2.
\]
特别地，\(q=9\) 是该类“阶塌缩”（Hankel 秩下降）在已核验窗内首次显露的最小 \(q\)。
该固定亏格现象表明：高阶 Rényi 碰撞矩并不必然引入同阶数量的新可观测谱模态；相反，可观测子空间可在离散 \(q\) 上发生代数消隐，从而出现可审计的“复杂度反直觉下降”。
\end{proposition}
\begin{proof}
次数与显式表达式见定理 \ref{thm:pom-moment-resonance}。
不可约性可由模素数不可约证书审计：例如表 \ref{tab:fold_collision_observable_minpoly_galois_certificate_q9_13} 的 \(p_{\mathrm{irr}}\) 列给出素数 \(p\) 使得 \(P_q(\lambda)\bmod p\) 仍不可约；从而 \(P_q\) 在 \(\QQ[\lambda]\) 上不可约，进而在 \(\ZZ[\lambda]\) 上不可约（Gauss 引理）。
阶塌缩等式只是把 \(2\lfloor q/2\rfloor+1\) 与上述次数直接相减。证毕。
\end{proof}

\begin{proposition}[共振窗 \texorpdfstring{$q=9,\dots,17$}{q=9..17} 的不可约性与模素数证书]\label{prop:pom-resonance-window-minpoly-irreducible-q9-17}
对共振窗 \(q\in\{9,10,\dots,17\}\)，由表 \ref{tab:fold_collision_moment_recursions_9_17} 的最小整系数递推系数所诱导的特征多项式 \(P_q(\lambda)\in\ZZ[\lambda]\) 在 \(\QQ[\lambda]\) 上不可约。
更精确地，可取素数证书
\[
(p_9,p_{10},p_{11},p_{12},p_{13},p_{14},p_{15},p_{16},p_{17})
=(11,17,37,29,29,37,17,239,31),
\]
使得对每个 \(q\in\{9,\dots,17\}\)，模 \(p_q\) 约化 \(P_q(\lambda)\bmod p_q\) 在 \(\FF_{p_q}[\lambda]\) 上不可约。
\end{proposition}
\begin{proof}
由表 \ref{tab:fold_collision_moment_recursions_9_17} 的递推
\(
S(m)=\sum_{i=1}^{d_q}c_i S(m-i)
\)
可定义首一多项式
\[
P_q(\lambda):=\lambda^{d_q}-\sum_{i=1}^{d_q}c_i\,\lambda^{d_q-i}\in\ZZ[\lambda].
\]
对每个 \(q\in\{9,\dots,17\}\)，选取上式所列素数 \(p_q\) 并在 \(\FF_{p_q}[\lambda]\) 中检验 \(P_q\bmod p_q\) 不可约，即得 \(P_q\) 在 \(\FF_{p_q}\) 上不可约。
由于 \(P_q\) 为本原首一整数多项式，若其在 \(\QQ[\lambda]\) 可约，则在 \(\ZZ[\lambda]\) 可约，从而在任意素数模约化下亦可约；
因此 \(P_q\bmod p_q\) 的不可约性推出 \(P_q\) 在 \(\QQ[\lambda]\) 上不可约。证毕。
\end{proof}

\begin{conjecture}[全对称伽罗瓦群与素数谱随机性]\label{conj:pom-resonance-minpoly-galois-sd}
对所有足够大的 \(q\)，可观测最小多项式 \(P_q(\lambda)\) 在 \(\QQ[\lambda]\) 上不可约且其伽罗瓦群满足
\[
\mathrm{Gal}(P_q/\QQ)\cong S_{\deg P_q}.
\]
在该假设下，由 Čebotarev 密度定理可预期：当素数 \(p\) 变化时，\(P_q(\lambda)\bmod p\) 的分解型在 \(S_{\deg P_q}\) 的共轭类上渐近等分布；
相应的模 \(p\) 递推动力学在素数谱意义下呈现接近最大混合的分解指纹。
\end{conjecture}

\begin{proposition}[判别式黑名单与分歧素数集：局部谱分歧仅可能发生在 \texorpdfstring{$p\mid \mathrm{Disc}(P_q)$}{p|Disc(Pq)}（\(q=9,10,11,13\)，可审计）]\label{prop:pom-observable-minpoly-discriminant-ramification-q9-13}
对 \(q\in\{9,10,11,13\}\)，记 \(\Delta_q^{\mathrm{disc}}\) 为 \(\mathrm{Disc}(P_q)\) 的平方自由部分（含符号），并定义伴随二次域
\[
K_q^{\mathrm{disc}}:=\QQ\!\left(\sqrt{\Delta_q^{\mathrm{disc}}}\right).
\]
令
\(
\mathcal{B}_q:=\{p:\ p\ \text{为素数且}\ p\mid \mathrm{Disc}(P_q)\}
\)
为分歧素数集。
则：
\begin{enumerate}
  \item \(\mathrm{Disc}(P_q)\) 的素因子分解由式 \eqref{eq:fold_collision_observable_minpoly_discriminants_q9_13} 给出；因此 \(\Delta_q^{\mathrm{disc}}\) 非平凡，且诱导二次域 \(K_q^{\mathrm{disc}}\)。其 \(\Delta_q^{\mathrm{disc}}\) 与 \(\mathcal{B}_q\) 的显式列表见表 \ref{tab:fold_collision_observable_minpoly_ramified_primes_q9_13}。
  \item 对任意素数 \(p\notin\mathcal{B}_q\)，多项式 \(P_q(\lambda)\bmod p\) 在 \(\FF_p[\lambda]\) 上必为平方自由；因而所有“\(\mathrm{repeat?}\)”型局部谱分歧（重根/Jordan 调制）只能发生在有限黑名单 \(\mathcal{B}_q\) 上。
\end{enumerate}
\end{proposition}
\begin{proof}
第 1 点由脚本 \path{scripts/exp_fold_collision_observable_minpoly_galois_audit_q9_13.py} 对定理 \ref{thm:pom-moment-resonance} 中 \(P_q\) 的判别式做符号计算并整因子分解得到，可复算输出即式 \eqref{eq:fold_collision_observable_minpoly_discriminants_q9_13}。
第 2 点是判别式的标准性质：对 \(\mathrm{char}(\FF_p)\nmid \mathrm{Disc}(P_q)\)（等价于 \(p\notin\mathcal{B}_q\)），有
\(\gcd(P_q,\partial_\lambda P_q)=1\) 于 \(\FF_p[\lambda]\)，故 \(P_q\bmod p\) 无重根，从而平方自由。
结合注 \ref{rem:pom-a2-embedding-local-ramification} 的局部判据，可把“异常素数集”在该窗口内完全操作化为 \(\mathcal{B}_q\)。证毕。
\end{proof}
\input{sections/generated/eq_fold_collision_observable_minpoly_discriminants_q9_13}
\input{sections/generated/tab_fold_collision_observable_minpoly_ramified_primes_q9_13}

\begin{proposition}[高阶递推的最大伽罗瓦群：\texorpdfstring{$\mathrm{Gal}(P_q/\QQ)\cong S_d$}{Gal(Pq/Q)=Sd}（\(q=9,10,11,13\)，可审计）]\label{prop:pom-observable-minpoly-galois-sd-q9-13}
设 \(d=\deg P_q\)。
则对 \(q\in\{9,10,11,13\}\) 有
\[
\mathrm{Gal}(P_9/\QQ)\cong S_7,\qquad
\mathrm{Gal}(P_{10}/\QQ)\cong S_9,\qquad
\mathrm{Gal}(P_{11}/\QQ)\cong S_9,\qquad
\mathrm{Gal}(P_{13}/\QQ)\cong S_{11}.
\]
\end{proposition}
\begin{proof}
对每个 \(q\in\{9,10,11,13\}\)，表 \ref{tab:fold_collision_observable_minpoly_galois_certificate_q9_13} 给出三个不分歧素数证书：
\begin{itemize}
  \item 存在素数 \(p_{\mathrm{irr}}\notin\mathcal{B}_q\) 使得 \(P_q\bmod p_{\mathrm{irr}}\) 不可约，于是 Frobenius 含 \(d\)-循环，且 \(P_q\) 在 \(\QQ\) 上不可约（从而 \(G_q:=\mathrm{Gal}(P_q/\QQ)\) 传递）；
  \item 存在素数 \(p_{(d-1,1)}\notin\mathcal{B}_q\) 使得分解型为 \((d-1,1)\)，故 \(G_q\) 含 \((d-1)\)-循环，从而 \(G_q\) 为原始群；
  \item 存在素数 \(p_{\mathrm{J}}\notin\mathcal{B}_q\) 使得分解型包含一条素数长度循环：对 \(d=7\) 含 \(3\)-循环、对 \(d=9\) 含 \(5\)-循环、对 \(d=11\) 含 \(7\)-循环；且该素数长度 \(\le d-3\)。
\end{itemize}
由 Jordan 定理，原始群 \(G_q\) 若含素数长度 \(p\)-循环（\(p\le d-3\)），则 \(A_d\subseteq G_q\)。
另一方面，由命题 \ref{prop:pom-observable-minpoly-discriminant-ramification-q9-13}，\(\mathrm{Disc}(P_q)\) 非平方（\(\Delta_q^{\mathrm{disc}}\neq 1\)），从而 \(G_q\not\subseteq A_d\)。
综上 \(G_q=S_d\)。证毕。
\end{proof}
\input{sections/generated/tab_fold_collision_observable_minpoly_galois_certificate_q9_13}

\begin{corollary}[Chebotarev 密度的可检验频率律：不可约密度为 \texorpdfstring{$1/d$}{1/d}]\label{cor:pom-observable-minpoly-chebotarev-density-q9-13}
在命题 \ref{prop:pom-observable-minpoly-galois-sd-q9-13} 的结论下，对去除有限分歧素数集 \(\mathcal{B}_q\) 后的素数集合，任意分裂型（等价于 \(S_d\) 中某个共轭类）出现的自然密度等于该共轭类在 \(S_d\) 中的比例。
特别地，\(P_q(\lambda)\bmod p\) 仍不可约（Frobenius 为 \(d\)-循环）的密度为 \(1/d\)，故
\[
\mathrm{dens}\{p:\ P_9\bmod p\ \text{不可约}\}=\frac17,\quad
\mathrm{dens}\{p:\ P_{10}\bmod p\ \text{不可约}\}=\frac19,
\]
\[
\mathrm{dens}\{p:\ P_{11}\bmod p\ \text{不可约}\}=\frac19,\quad
\mathrm{dens}\{p:\ P_{13}\bmod p\ \text{不可约}\}=\frac1{11}.
\]
\end{corollary}
\begin{proof}
由 Chebotarev 密度定理，未分歧素数的 Frobenius 共轭类在 \(G_q=S_d\) 中按共轭类大小等分布；而 \(d\)-循环在 \(S_d\) 中所占比例为 \((d-1)!/d!=1/d\)。证毕。
\end{proof}

\begin{corollary}[判别式二次特征给出 Frobenius 奇偶的闭式判据]\label{cor:pom-observable-minpoly-frob-parity-q9-13}
在 \(G_q=S_d\) 下，\(S_d\) 的唯一指数 \(2\) 正规子群为 \(A_d\)，其固定子域为 \(\QQ(\sqrt{\mathrm{Disc}(P_q)})=\QQ(\sqrt{\Delta_q^{\mathrm{disc}}})\)。
因此对任意不分歧素数 \(p\notin\mathcal{B}_q\)，Frobenius 置换的奇偶性由 Kronecker 符号严格给出：
\[
\mathrm{sgn}(\mathrm{Frob}_p)=\left(\frac{\Delta_q^{\mathrm{disc}}}{p}\right).
\]
特别地，对 \(q=9\)，由式 \eqref{eq:fold_collision_observable_minpoly_discriminants_q9_13} 可知 \(\Delta_9^{\mathrm{disc}}\) 为素数，从而奇偶判据退化为单一 Legendre 符号 \(\left(\frac{\Delta_9^{\mathrm{disc}}}{p}\right)\)。
\end{corollary}
\begin{proof}
这是判别式的标准刻画：\(\mathrm{Disc}(P_q)\) 的平方类对应 \(G_q\) 到 \(\{\pm 1\}\) 的符号表示（置换奇偶）所切出的二次子域；对不分歧素数，Frobenius 的像由二次互反律的 Kronecker 符号读出。证毕。
\end{proof}

\begin{proposition}[负实主导律与交替修正：\texorpdfstring{$q=9,10,11,13$}{q=9,10,11,13} 的有效 Ramanujan 比率]\label{prop:pom-observable-minpoly-negative-real-dominance-q9-13}
对 \(q\in\{9,10,11,13\}\)，记 \(r_q\) 为 \(P_q\) 的 Perron 根，并令
\[
\Lambda_q:=\max\{|\mu|:\ P_q(\mu)=0,\ \mu\neq r_q\}.
\]
则达到 \(\Lambda_q\) 的根为唯一负实根 \(\mu^-_q<0\)，从而 \(S_q(m)\) 的首个非主指数修正带严格交替相位 \(({-1})^m\)。
并且有效 Ramanujan 比率
\(
R_q:=|\mu^-_q|/\sqrt{r_q}
\)
在该窗口内满足式 \eqref{eq:fold_collision_observable_minpoly_negative_real_dominance_q9_13} 的可复算数值（均严格大于 \(1\)）。
\end{proposition}
\begin{proof}
由脚本 \path{scripts/exp_fold_collision_observable_minpoly_galois_audit_q9_13.py} 对 \(P_q\) 的全部根做高精度数值求解，并逐一比较模长得到：除 Perron 根外，最大模根为唯一负实根。
相应地，递推解的谱展开中该根贡献的符号为 \((\mu^-_q)^m=(-1)^m|\mu^-_q|^m\)，故出现交替修正；比率 \(R_q\) 的数值即由该根与 \(r_q\) 直接代入定义得到。证毕。
\end{proof}
\input{sections/generated/eq_fold_collision_observable_minpoly_negative_real_dominance_q9_13}

